% sci_build_flash_flow.tex — Vertical edit/build/flash/monitor cycle for the
% SCI/SICE tutorial "learner code workflow" slide (S2, P37).
% SCI/SICEチュートリアル「学習者コードの型とワークフロー」スライド用の
% 縦型 編集/ビルド/書込/モニタ 周回図。
%
% Not shared with the DXH workshop deck (build_flash_flow.tex, horizontal,
% 4 steps, no edit step) — kept as a separate file per that deck's request.
% DXHワークショップ資料が使う build_flash_flow.tex（横型・4ステップ・
% editステップなし）とは別ファイル（先方の指定により共有しない）。
\documentclass[border=6pt]{standalone}
\usepackage{tikz}
\usetikzlibrary{arrows.meta, positioning, calc}
\usepackage{luatexja}
\usepackage[match]{luatexja-fontspec}
% Cross-platform font: Hiragino (macOS) -> Yu Gothic (Windows) -> Noto (Linux)
% クロスプラットフォーム: ヒラギノ (macOS) -> 游ゴシック (Windows) -> Noto (Linux)
\usepackage{ifthen}
\newboolean{jfontset}\setboolean{jfontset}{false}
\IfFontExistsTF{Hiragino Kaku Gothic ProN}{%
  \setmainjfont{Hiragino Kaku Gothic ProN}%
  \setsansjfont{Hiragino Kaku Gothic ProN}%
  \setboolean{jfontset}{true}}{}
\ifthenelse{\boolean{jfontset}}{}{%
  \IfFontExistsTF{Yu Gothic}{%
    \setmainjfont{Yu Gothic}%
    \setsansjfont{Yu Gothic}%
    \setboolean{jfontset}{true}}{}}
\ifthenelse{\boolean{jfontset}}{}{%
  \IfFontExistsTF{Noto Sans CJK JP}{%
    \setmainjfont{Noto Sans CJK JP}%
    \setsansjfont{Noto Sans CJK JP}}{}}

\begin{document}

\definecolor{blue1}{RGB}{0, 100, 180}
\definecolor{orange1}{RGB}{200, 110, 0}
\definecolor{green1}{RGB}{50, 140, 50}
\definecolor{yellow1}{RGB}{180, 140, 0}
\definecolor{purple1}{RGB}{120, 80, 150}
\definecolor{gray1}{RGB}{120, 120, 120}

\begin{tikzpicture}[
    step/.style={
      rectangle, rounded corners=5pt, draw=#1, line width=1.5pt,
      fill=#1!8, minimum width=40mm, minimum height=10mm,
      font=\bfseries\normalsize, text=#1, align=center, inner sep=2pt
    },
    desc/.style={font=\scriptsize, text=gray1, align=left},
    arr/.style={-{Stealth[length=3.5mm, width=2.5mm]}, line width=1.8pt, #1},
    feedback/.style={-{Stealth[length=3mm, width=2mm]}, line width=1.1pt, gray1, dashed},
  ]

  % Steps, top to bottom
  % ステップ（上から下へ）
  \node[step=blue1]                 (s1) at (0, 0) {sf lesson\\switch N};
  \node[step=orange1, below=8mm of s1] (s2) {sf lesson\\edit};
  \node[step=green1,  below=8mm of s2] (s3) {sf lesson\\build};
  \node[step=yellow1, below=8mm of s3] (s4) {sf lesson\\flash};
  \node[step=purple1, below=8mm of s4] (s5) {sf monitor\\workshop};

  % Main flow arrows (right-angle: each is already a single vertical
  % segment since the nodes share the same x-coordinate)
  % 本流の矢印（ノードが同一x座標のため各区間は単純な垂直線）
  \draw[arr=blue1]   (s1) -- (s2);
  \draw[arr=orange1] (s2) -- (s3);
  \draw[arr=green1]  (s3) -- (s4);
  \draw[arr=yellow1] (s4) -- (s5);

  % Descriptions to the right of each box
  % 各ボックス右側の説明
  \node[desc, right=4mm of s1] {レッスンへ切替};
  \node[desc, right=4mm of s2] {user\_code.cpp\\を書く};
  \node[desc, right=4mm of s3] {ESP-IDF\\コンパイル};
  \node[desc, right=4mm of s4] {USB書込};
  \node[desc, right=4mm of s5] {波形/出力を確認};

  % Feedback loop: from the confirm step (s5) back up to edit (s2), for
  % iterating on the code. Routed on the LEFT side (west) so it does not
  % cross the description labels on the right, using right-angle routing
  % (horizontal - vertical - horizontal) with >=5mm of straight stem before
  % the arrowhead enters s2.
  % 確認(s5)から編集(s2)へ戻るフィードバック。右側の説明ラベルと交差しない
  % よう左(west)側を通す。直角ルーティング(水平-垂直-水平)、矢頭手前に
  % 5mm以上の直線区間を確保。
  % NOTE: (a -| b) is (b.x, a.y). The vertical run must sit at fbx's x, so
  % the corner at s2's level is (s2.west -| fbx), not (fbx -| s2.west).
  % 注意: (a -| b) は「b の x、a の y」。縦線は fbx の x に置くので、
  % s2 の高さの角は (s2.west -| fbx)。逆に書くと s2 の下から入ってしまう。
  \coordinate (fbx) at ($(s5.west) + (-26mm, 0)$);
  \coordinate (fbtop) at (s2.west -| fbx);
  \draw[feedback] (s5.west) -- (fbx) -- (fbtop) -- (s2.west);
  \node[font=\scriptsize\itshape, text=gray1, align=right, anchor=east, fill=white, inner sep=1pt]
    at ($(fbx)!0.5!(fbtop) + (-2mm, 0)$) {直して\\再挑戦};

\end{tikzpicture}
\end{document}
